
\section*{Conclusion générale}

Ce projet de simulation multi-agents, intitulé \textit{Gardien du parc}, s’est révélé être une expérience riche, transversale et particulièrement formatrice. À travers la conception, l’implémentation et l’analyse d’un système complexe, nous avons pu mobiliser et approfondir un large éventail de compétences acquises au cours de notre formation universitaire.

Sur le plan technique, la réalisation du simulateur nous a permis de consolider notre maîtrise du langage Java et de ses paradigmes fondamentaux tels que l’orientation objet, la programmation événementielle et la gestion des interfaces graphiques via Swing. L’utilisation rigoureuse du modèle MVC, le recours aux design patterns , et l’intégration de composants avancés témoignent de notre capacité à structurer un projet logiciel selon les standards professionnels.

La simulation d’agents autonomes évoluant dans un environnement contraint a nécessité la mise en œuvre de mécanismes de perception, de décision et d’action. Les comportements des gardiens et des intrus, que nous avons modélisés à travers différentes stratégies de déplacement, ont offert un terrain d’expérimentation concret pour explorer les bases de l’intelligence artificielle distribuée. Les scénarios testés nous ont permis de valider nos algorithmes, de mesurer leur efficacité, et de tirer des enseignements sur la coordination entre agents.

En parallèle, la couche graphique a donné une dimension interactive et visuelle à notre projet, renforçant sa lisibilité et sa pertinence. L’affichage en temps réel des déplacements, du champ de vision, des captures et des obstacles a permis une meilleure compréhension du système, tant pour les concepteurs que pour les observateurs extérieurs. L’ajout de visualisations statistiques grâce à JFreeChart a transformé notre application en un véritable outil d’analyse, capable de produire des indicateurs mesurables et exploitables.

Au-delà de l’aspect technique, ce projet nous a confrontés à la réalité du travail collaboratif. La gestion de versions avec GitHub, la répartition des responsabilités, les revues de code croisées et la rédaction conjointe du rapport via Overleaf ont structuré notre démarche et renforcé notre esprit d’équipe. Ces pratiques sont représentatives de celles utilisées dans le monde professionnel, et nous ont permis d’acquérir des compétences précieuses en gestion de projet, en communication et en organisation.

Nous avons également développé une autonomie dans le choix des outils et des extensions. Face aux contraintes initiales, nous avons su élargir notre périmètre technologique de manière cohérente et justifiée. Cela s’est traduit par une meilleure adaptabilité, une capacité à résoudre des problèmes imprévus, et une amélioration constante de la qualité du code et de l’expérience utilisateur.
Ce projet n’a pas été sans défis. Même si nous avons utilisé un seul thread pour piloter l’ensemble des agents, la gestion fluide des déplacements, la réactivité de l’affichage et la modularité du code ont exigé rigueur et persévérance. Chaque difficulté surmontée a renforcé notre compréhension des enjeux du développement logiciel et de la complexité des systèmes interactifs.

En conclusion, ce travail nous a permis de franchir une étape importante dans notre parcours. Il nous a donné l’occasion d’appliquer nos connaissances de manière concrète, de collaborer efficacement en équipe, et de réaliser une application complète, fiable et évolutive. Le projet Gardien du parc constitue une vraie synthèse de nos apprentissages en génie logiciel, et ouvre la voie vers des domaines passionnants comme les simulations interactives, l'intelligence artificielle distribuée, ou encore le développement d'applications logicielles à grande échelle.

Nous gardons de cette expérience un souvenir formateur et motivant, qui marque une étape importante dans notre progression vers les métiers de l'informatique : développeur, concepteur logiciel, spécialiste en simulation, architecte logiciel, et bien d'autres.



\end{document}